%----------------------------------------------------------------------------------------
% REKENKUNDIGE EN MEETKUNDIGE RIJEN. Alle opgaven zijn eigen werk.
%
% Opbouw: Start, leerstof, en op een nieuwe pagina de Opdrachten, Koppel en Test jezelf.
% De oefeningen staan in oefeningen/02-rekenkundige-en-meetkundige-rijen.yaml.
% Labels van dit hoofdstuk beginnen met h2-.
%----------------------------------------------------------------------------------------

\mychapter{2}{Rekenkundige en meetkundige rijen}

\lkr{Leerplan B+S'': LPD 44, zonder de somformules (die staan in hoofdstuk 3). Twee rode draden: de sprongenfiguur ($n-1$ sprongen, zowel bij $+v$ als bij $\cdot\, q$) en het onderscheid tussen rekenkundig en meetkundig, dat in de gemengde Opdrachten steeds terugkomt. De formule voor exponentiële groei wordt afgeleid uit een getallenvoorbeeld: laat leerlingen die afleiding zelf reconstrueren.}

In hoofdstuk 1 zag je dat een rij allerlei vormen kan aannemen. Twee soorten komen zo vaak voor dat ze een eigen naam krijgen. Bij een \textit{rekenkundige} rij komt er telkens hetzelfde getal bij, zoals bij een stapel stoelen die per stoel even veel hoger wordt. Bij een \textit{meetkundige} rij vermenigvuldig je telkens met hetzelfde getal, zoals bij een bedrag dat elk jaar met hetzelfde percentage groeit. Voor beide soorten vind je in dit hoofdstuk een expliciet voorschrift, en je ontdekt dat ze de discrete versie zijn van twee functies die je al kent: de eerstegraadsfunctie en de exponentiële functie.

\begin{oefeningen}[Start]
\oef{h2-start-1}

\oef{h2-verken}
\end{oefeningen}

\lkr{Laat oefening \ref{oef:h2-verken} maken vóór de uitleg, zonder formule. Vergelijk daarna de antwoorden op b: wie rekent $25\cdot 5$ en wie $24\cdot 5$, en waarom?}

%========================================================================================
\section{Rekenkundige rijen}

\subsection{Definitie en voorschriften}

In de verkenning werd de stapel bij elke stoel 5 cm hoger. De hoogtes $80,\ 85,\ 90,\ \ldots$ vormen een rij waarin er telkens hetzelfde getal bijkomt. Zo'n rij heeft een eenvoudig recursief voorschrift, en daaruit volgt meteen een expliciet voorschrift.

\needspace{20\baselineskip}
\begin{theorie}
Een rij is \emph{rekenkundig} als elke term de vorige term plus hetzelfde getal is. Dat getal is het \emph{verschil} $v$:

\qquad $\formule{u_n = u_{n-1} + v}$ \quad voor $n \geqslant 2$, met $u_1$ gegeven

\begin{stelling}{Expliciet voorschrift van een rekenkundige rij}
\begin{minipage}[t]{.4\linewidth}
\vspace{0pt}
\qquad $\formule{u_n = u_1 + (n-1)\,v}$
\end{minipage}\hfill
\begin{minipage}[t]{.56\linewidth}
\vspace{0pt}
\sprongfiguur{+v}
\end{minipage}
\end{stelling}

\begin{bewijs}
Van $u_1$ naar $u_n$ zet je $n - 1$ sprongen, en bij elke sprong komt er $v$ bij. Dus
\bewijsregel{u_n = u_1 + \underbrace{v + v + \cdots + v}_{n-1} = u_1 + (n-1)\,v .}
\end{bewijs}
\end{theorie}

De valkuil zit in het aantal sprongen. Voor 25 stoelen tel je geen 25 keer 5 cm bij de eerste stoel, want de eerste stoel zelf is geen sprong: het zijn er 24.

\voorbeeld{}{Bepaal het expliciet voorschrift van $7,\ 11,\ 15,\ \ldots$ en bereken $u_{30}$.\label{vb:h2-term}}{%
\uitleg{Voor het voorschrift hebben we $u_1$ en $v$ nodig. De eerste term staat er, het verschil halen we uit twee opeenvolgende termen.}
\begin{stappen}
v = 11 - 7 = 4 & Het verschil is telkens 4. We controleren met het volgende paar: $15 - 11 = 4$. \\
u_n = 7 + (n-1) \cdot 4 = 4n + 3 & We vullen $u_1 = 7$ en $v = 4$ in, en vereenvoudigen. \\
u_{30} = 4 \cdot 30 + 3 = 123 & We vullen $n = 30$ in. Ter controle: $7 + 29 \cdot 4 = 123$, met 29 sprongen.
\end{stappen}}

\nujij{Bepaal het expliciet voorschrift van $12,\ 9,\ 6,\ \ldots$ en bereken $u_{25}$.}{$v = -3$, $u_n = 15 - 3n$, $u_{25} = -60$.}

\thuis{\blokkop{Half uitgewerkt}\quad Rij $3;\ 3{,}5;\ 4;\ \ldots$ Bereken $u_{41}$. Schrijf over in je schrift en vul aan: \ $v = \vak$, \ $u_{41} = 3 + \vak \cdot \vak = \vak$. \quad {\footnotesize(Antwoord: $v = 0{,}5$ en $u_{41} = 3 + 40 \cdot 0{,}5 = 23$.)}}

\begin{oefeningen}
\oef{h2-k1}
\end{oefeningen}

\subsection{Term of geen term}

Net zoals in hoofdstuk 1 kun je nagaan of een getal in de rij voorkomt: stel de algemene term gelijk aan dat getal, en los op naar $n$. Bij een rekenkundige rij is dat een eerstegraadsvergelijking.

\voorbeeld{}{Is 203 een term van $7,\ 11,\ 15,\ \ldots$? En 200?\label{vb:h2-isterm}}{%
\uitleg{Uit voorbeeld \ref{vb:h2-term} kennen we het voorschrift. We stellen het gelijk aan het getal, en kijken of $n$ een volgnummer is.}
\begin{stappen}
7 + (n-1) \cdot 4 = 203 \ \Rightarrow \ n = 50 & $n - 1 = 49$, dus $n = 50$. Dat is een volgnummer: 203 is de 50ste term. \\
7 + (n-1) \cdot 4 = 200 \ \Rightarrow \ n = 49{,}25 & Dat is geen natuurlijk getal, dus 200 is geen term. Je ziet het ook zonder rekenen: $200 - 7 = 193$ is geen veelvoud van 4.
\end{stappen}}

\begin{werkwijze}
Los $u_n = x$ op. Is $n$ een natuurlijk getal (minstens 1), dan is $x$ de $n$-de term. Anders is $x$ geen term.
\end{werkwijze}

\nujij{Is 10 een term van $100,\ 94,\ 88,\ \ldots$? En 1?}{10 is de 16de term; bij 1 is $n = 17{,}5$, dus geen term.}

\begin{oefeningen}
\oef{h2-k2}
\end{oefeningen}

\subsection{De rij vinden uit twee termen}

Soms ken je niet de eerste term, maar twee termen ergens anders in de rij. Dezelfde redenering met sprongen werkt ook tussen die twee termen: van $u_m$ naar $u_n$ zijn het $n - m$ sprongen.

\begin{theorie}
\begin{stelling}{Twee termen van een rekenkundige rij}
\qquad $\formule{u_n = u_m + (n-m)\,v}$
\end{stelling}
\end{theorie}

\voorbeeld{}{$u_3 = 11$ en $u_9 = 35$. Bepaal $u_1$ en $v$ op twee manieren: met een stelsel en met sprongen.\label{vb:h2-tweetermen}}{%
\uitleg{Een stelsel werkt altijd, maar is wat meer schrijfwerk. Met sprongen gaat het sneller. We vergelijken de twee.}

\begin{minipage}[t]{.47\linewidth}
\textbf{Stelsel.} \ We schrijven beide termen met het voorschrift: $u_1 + 2v = 11$ en $u_1 + 8v = 35$.\\
De tweede vergelijking min de eerste geeft $6v = 24$, dus $v = 4$.\\
Invullen: $u_1 = 11 - 2 \cdot 4 = 3$.
\end{minipage}\hfill
\begin{minipage}[t]{.47\linewidth}
\textbf{Sprongen.} \ Van $u_3$ naar $u_9$ zijn het 6 sprongen.\\
Ze overbruggen $35 - 11 = 24$, dus $6v = 24$ en $v = 4$.\\
Twee sprongen terug: $u_1 = 11 - 2 \cdot 4 = 3$.
\end{minipage}

\vspace{1mm}
In het stelsel zit dezelfde berekening: het aftrekken van de vergelijkingen telt net de sprongen tussen $u_3$ en $u_9$.}

\nujij{$u_4 = -2$ en $u_{10} = 16$. Bepaal $u_1$ en $v$.}{$v = 3$ en $u_1 = -11$.}

\begin{oefeningen}
\oef{h2-k3}
\end{oefeningen}

\subsection{Rekenkundig midden}

In een rekenkundige rij ligt elke term precies in het midden tussen zijn twee buren: hij is een sprong verder dan de vorige, en een sprong minder ver dan de volgende.

\begin{theorie}
\begin{stelling}{Rekenkundig midden}
In een rekenkundige rij is elke term (behalve de eerste) het gemiddelde van de twee termen ernaast:

\vspace{1mm}
\qquad $\formule{u_n = \dfrac{u_{n-1} + u_{n+1}}{2}}$
\end{stelling}

\begin{bewijs}
Een sprong terug en een sprong vooruit geven $u_{n-1} = u_n - v$ en $u_{n+1} = u_n + v$. Tel ze op: $-v$ en $+v$ vallen weg, zodat
\bewijsregel{u_{n-1} + u_{n+1} = 2u_n , \quad\text{dus}\quad u_n = \dfrac{u_{n-1} + u_{n+1}}{2} .}
\end{bewijs}
\end{theorie}

Omgekeerd zijn drie getallen opeenvolgende termen van een rekenkundige rij als het middelste het gemiddelde is van de twee andere. Daarmee kun je een onbekende bepalen.

\voorbeeld{}{Voor welke $x$ zijn $2x$, $x + 5$ en $3x + 1$ drie opeenvolgende termen van een rekenkundige rij?\label{vb:h2-midden}}{%
\uitleg{Het middelste getal moet het gemiddelde zijn van de twee andere. Dat geeft een vergelijking in $x$.}
\begin{stappen}
x + 5 = \dfrac{2x + (3x + 1)}{2} & We schrijven het rekenkundig midden uit. \\[1.5ex]
2x + 10 = 5x + 1 \ \Rightarrow \ x = 3 & Beide leden maal 2, en oplossen. \\
6,\ 8,\ 10 & We controleren door in te vullen: het verschil is telkens 2.
\end{stappen}}

\begin{oefeningen}
\oef{h2-k4}
\end{oefeningen}

\subsection{Grafiek}

Werk het voorschrift uit, en je ziet een bekende vorm: $u_n = u_1 + (n-1)\,v = v\,n + (u_1 - v)$. Dat is het voorschrift van een eerstegraadsfunctie, met $n$ in plaats van $x$. De punten van een rekenkundige rij liggen dus op een rechte met richtingscoëfficiënt $v$. Een stap van 1 naar rechts geeft een stap van $v$ omhoog: bij $v > 0$ stijgt de rij, bij $v < 0$ daalt ze. Alleen de punten horen bij de rij, de rechte niet.

\figuren{\begin{tikzpicture}[x=.55cm,y=.2cm,font=\scriptsize]
  \draw[black!14,very thin,xstep=1,ystep=2] (0,0) grid (6.4,18);
  \draw[->,grijs] (0,0) -- (6.8,0) node[right,text=black] {$n$};
  \draw[->,grijs] (0,0) -- (0,19) node[above,text=black] {$u_n$};
  \foreach \t in {1,...,6}{\node[below,inner sep=1.5pt] at (\t,0) {$\t$};}
  \foreach \t in {4,8,12,16}{\node[left,inner sep=1.5pt] at (0,\t) {$\t$};}
  \draw[grijs,dashed] (.4,.2) -- (6.3,17.9);
  \draw[accent,thick] (3,8) -- node[below=-1pt,text=black] {$1$} (4,8) -- node[right=-1pt,text=black] {$v = 3$} (4,11);
  \foreach \x/\y in {1/2,2/5,3/8,4/11,5/14,6/17} \fill[accent] (\x,\y) circle (1.9pt);
  \node[above] at (3.2,19) {$2,\ 5,\ 8,\ \ldots$};
\end{tikzpicture}\qquad\qquad
\begin{tikzpicture}[x=.55cm,y=.2cm,font=\scriptsize]
  \draw[black!14,very thin,xstep=1,ystep=2] (0,0) grid (6.4,18);
  \draw[->,grijs] (0,0) -- (6.8,0) node[right,text=black] {$n$};
  \draw[->,grijs] (0,0) -- (0,19) node[above,text=black] {$u_n$};
  \foreach \t in {1,...,6}{\node[below,inner sep=1.5pt] at (\t,0) {$\t$};}
  \foreach \t in {4,8,12,16}{\node[left,inner sep=1.5pt] at (0,\t) {$\t$};}
  \draw[grijs,dashed] (.4,17.2) -- (6.4,5.2);
  \draw[accent,thick] (3,12) -- node[above=-1pt,text=black] {$1$} (4,12) -- node[right=-1pt,text=black] {$v = -2$} (4,10);
  \foreach \x/\y in {1/16,2/14,3/12,4/10,5/8,6/6} \fill[accent] (\x,\y) circle (1.9pt);
  \node[above] at (3.2,19) {$16,\ 14,\ 12,\ \ldots$};
\end{tikzpicture}}

\subsection{Lineaire groei}

Rekenkundige rijen beschrijven alles wat per tijdseenheid met hetzelfde bedrag toeneemt of afneemt. Neem een spaarpot met \eur 50, waar je elke week \eur 5 in stopt. Na 1 week zit er $50 + 5$ in, na 2 weken $50 + 2 \cdot 5$, en na $n$ weken $50 + n \cdot 5$. Omdat we hier beginnen bij de beginwaarde, voor er iets gebeurd is, start de rij bij $u_0 = 50$. Na $n$ weken zijn er dan precies $n$ sprongen gezet, en niet $n - 1$.

\begin{theorie}
\emph{Lineaire groei}: neemt een grootheid per tijdseenheid toe met hetzelfde bedrag $v$ (of af, bij $v < 0$), en is $u_0$ de beginwaarde, dan is ze na $n$ tijdseenheden

\vspace{1mm}
\qquad $\formule{u_n = u_0 + n\,v}$
\end{theorie}

\voorbeeld{}{Er staat \eur 350 op een rekening. Elke maand komt er \eur 25 bij. Na hoeveel maanden staat er minstens \eur 1000 op?\label{vb:h2-groei}}{%
\uitleg{Er komt elke maand hetzelfde bedrag bij, dus dit is lineaire groei. De \eur 350 is het bedrag bij de start, voor er een maand voorbij is: dat is $u_0$.}
\begin{stappen}
u_n = 350 + 25n & Met $u_0 = 350$ en $v = 25$ is dit het bedrag na $n$ maanden. \\
350 + 25n \geqslant 1000 & We zoeken wanneer het bedrag minstens 1000 is. \\
n \geqslant 26 & $25n \geqslant 650$ geeft $n \geqslant 26$. Na 26 maanden staat er precies \eur 1000 op.
\end{stappen}}

\lkr{Vergelijk met de stoelen: daar start de rij bij $u_1$, hier bij $u_0$. Laat leerlingen uitleggen waarom het voorschrift dan $n$ bevat en niet $n-1$.}

\nujij{Een waterton bevat 1200 liter en verliest 15 liter per uur. Geef een voorschrift voor de inhoud na $n$ uur. Na hoeveel uur is de ton leeg?}{$u_n = 1200 - 15n$; leeg na 80 uur.}

\begin{oefeningen}
\oef{h2-k5}
\end{oefeningen}

%========================================================================================
\section{Meetkundige rijen}

Bij een rekenkundige rij tel je telkens hetzelfde getal op. Wat gebeurt er als je in plaats daarvan telkens met hetzelfde getal vermenigvuldigt? Onderzoek het eerst zelf.

\begin{oefeningen}
\oef{h2-mverken}
\end{oefeningen}

\lkr{Laat de leerlingen de vergelijking maken met de stoelen: wat blijft hetzelfde (het aantal sprongen) en wat verandert (optellen wordt vermenigvuldigen)?}

\subsection{Definitie en voorschriften}

Het gerucht verspreidde zich als $1,\ 3,\ 9,\ 27,\ \ldots$: elke dag maal 3. Na een paar weken zijn de aantallen enorm, terwijl een rekenkundige rij in dezelfde tijd maar traag groeit. De sprongen zijn er nog altijd, maar nu is elke sprong een vermenigvuldiging.

\needspace{20\baselineskip}
\begin{theorie}
Een rij is \emph{meetkundig} als elke term de vorige term maal hetzelfde getal is. Dat getal is het \emph{quotiënt} $q$, met $q \neq 0$ en $u_1 \neq 0$:

\qquad $\formule{u_n = u_{n-1} \cdot q}$ \quad voor $n \geqslant 2$, met $u_1$ gegeven

\begin{stelling}{Expliciet voorschrift van een meetkundige rij}
\begin{minipage}[t]{.4\linewidth}
\vspace{0pt}
\qquad $\formule{u_n = u_1 \cdot q^{\,n-1}}$
\end{minipage}\hfill
\begin{minipage}[t]{.56\linewidth}
\vspace{0pt}
\sprongfiguur{\cdot\, q}
\end{minipage}
\end{stelling}

\begin{bewijs}
Van $u_1$ naar $u_n$ zet je $n - 1$ sprongen, en bij elke sprong vermenigvuldig je met $q$. Dus
\bewijsregel{u_n = u_1 \cdot \underbrace{q \cdot q \cdots q}_{n-1} = u_1 \cdot q^{\,n-1} .}
\end{bewijs}
\end{theorie}

Het quotiënt vind je door een term te delen door de term ervoor. Een rij met een 0 erin kan dus niet meetkundig zijn.

\voorbeeld{}{Bepaal het expliciet voorschrift van $3,\ 6,\ 12,\ \ldots$ en bereken $u_{10}$.\label{vb:h2-mterm}}{%
\uitleg{Elke term is het dubbele van de vorige, dus de rij is meetkundig. Voor het voorschrift hebben we $u_1$ en $q$ nodig.}
\begin{stappen}
q = \dfrac{6}{3} = 2 & We delen een term door de vorige. Controle met het volgende paar: $\dfrac{12}{6} = 2$. \\[1.5ex]
u_n = 3 \cdot 2^{\,n-1} & We vullen $u_1 = 3$ en $q = 2$ in. Dit mag je niet samenvoegen tot $6^{\,n-1}$: de macht hoort alleen bij de 2. \\
u_{10} = 3 \cdot 2^{9} = 1536 & We vullen $n = 10$ in: van $u_1$ naar $u_{10}$ zijn het 9 sprongen.
\end{stappen}}

\nujij{Bepaal het expliciet voorschrift van $4,\ {-8},\ 16,\ \ldots$ en bereken $u_7$.}{$q = -2$, $u_n = 4 \cdot (-2)^{n-1}$, $u_7 = 256$.}

\begin{oefeningen}
\oef{h2-k6}
\end{oefeningen}

\subsection{Term of geen term}

Ook hier stel je de algemene term gelijk aan het getal. Nu staat $n$ in de exponent: dat is een exponentiële vergelijking, die je oplost met een logaritme of door het getal als macht te herkennen.

\voorbeeld{}{Is 768 een term van $3,\ 6,\ 12,\ \ldots$? En 1000?\label{vb:h2-misterm}}{%
\uitleg{Het voorschrift is $u_n = 3 \cdot 2^{\,n-1}$ (voorbeeld \ref{vb:h2-mterm}). We maken de macht eerst alleen, en zoeken dan de exponent.}
\begin{stappen}
3 \cdot 2^{\,n-1} = 768 \ \Rightarrow \ 2^{\,n-1} = 256 & We delen beide leden door 3. \\
n - 1 = 8 \ \Rightarrow \ n = 9 & Omdat $2^8 = 256$, is $n = 9$: 768 is de 9de term. \\
3 \cdot 2^{\,n-1} = 1000 \ \Rightarrow \ 2^{\,n-1} = 333{,}33\ldots & Hetzelfde voor 1000. \\
n - 1 = \log_2 333{,}33\ldots \approx 8{,}38 & Geen natuurlijk getal, dus 1000 is geen term. Het ligt tussen $u_9 = 768$ en $u_{10} = 1536$.
\end{stappen}}

\begin{werkwijze}
Los $u_n = x$ op: maak de macht alleen, en bepaal de exponent, zo nodig met een logaritme. Is $n$ een natuurlijk getal (minstens 1), dan is $x$ de $n$-de term.
\end{werkwijze}

\nujij{Is 405 een term van $5,\ 15,\ 45,\ \ldots$? En 500?}{405 is de 5de term; bij 500 is $3^{n-1} = 100$ en $n - 1 \approx 4{,}19$, dus geen term.}

\begin{oefeningen}
\oef{h2-k7}
\end{oefeningen}

\subsection{De rij vinden uit twee termen}

Ken je twee termen, dan tel je weer de sprongen ertussen. Alleen wordt elke sprong nu een factor $q$.

\begin{theorie}
\begin{stelling}{Twee termen van een meetkundige rij}
\qquad $\formule{u_n = u_m \cdot q^{\,n-m}}$
\end{stelling}
\end{theorie}

Let op bij een even aantal sprongen. Uit $q^2 = 4$ volgt $q = 2$ of $q = -2$, en beide rijen voldoen.

\voorbeeld{}{$u_3 = 12$ en $u_5 = 48$. Bepaal $u_1$ en $q$.\label{vb:h2-mtwee}}{%
\uitleg{Van $u_3$ naar $u_5$ zijn het 2 sprongen. Dat geeft een vergelijking in $q^2$, en die heeft twee oplossingen.}
\begin{stappen}
12 \cdot q^2 = 48 \ \Rightarrow \ q^2 = 4 & Twee sprongen vooruit vanaf $u_3$. \\
q = 2 \quad\text{of}\quad q = -2 & Een even macht: er zijn twee oplossingen. \\[1.5ex]
u_1 = \dfrac{12}{q^2} = 3 & Twee sprongen terug vanaf $u_3$. Omdat $q^2 = 4$ in beide gevallen, is $u_1 = 3$ in beide gevallen. \\[1.5ex]
3,\ 6,\ 12,\ 24,\ 48 \quad\text{of}\quad 3,\ {-6},\ 12,\ {-24},\ 48 & We controleren: beide rijen hebben $u_3 = 12$ en $u_5 = 48$.
\end{stappen}}

\nujij{$u_2 = 4$ en $u_4 = 36$. Bepaal $u_1$ en $q$.}{$q = 3$ en $u_1 = \tfrac43$, of $q = -3$ en $u_1 = -\tfrac43$.}

\begin{oefeningen}
\oef{h2-k8}
\end{oefeningen}

\subsection{Meetkundig midden}

Bij een rekenkundige rij is elke term het gemiddelde van zijn buren. Bij een meetkundige rij speelt vermenigvuldigen dezelfde rol als optellen: een sprong terug is delen door $q$, een sprong vooruit is vermenigvuldigen met $q$.

\begin{theorie}
\begin{stelling}{Meetkundig midden}
In een meetkundige rij is het kwadraat van elke term (behalve de eerste) het product van de twee termen ernaast:

\vspace{1mm}
\qquad $\formule{u_n^{\,2} = u_{n-1} \cdot u_{n+1}}$
\end{stelling}

\begin{bewijs}
Een sprong terug en een sprong vooruit geven $u_{n-1} = \dfrac{u_n}{q}$ en $u_{n+1} = u_n \cdot q$. Vermenigvuldig ze: $q$ valt weg, zodat
\bewijsregel{u_{n-1} \cdot u_{n+1} = \dfrac{u_n}{q} \cdot u_n \cdot q = u_n^{\,2} .}
\end{bewijs}
\end{theorie}

\voorbeeld{}{Voor welke $x$ zijn $x + 1$, $x + 5$ en $x + 17$ drie opeenvolgende termen van een meetkundige rij?\label{vb:h2-mmidden}}{%
\uitleg{Het kwadraat van het middelste getal moet het product zijn van de twee andere. Dat geeft een vergelijking in $x$.}
\begin{stappen}
(x + 5)^2 = (x + 1)(x + 17) & We schrijven het meetkundig midden uit. \\
x^2 + 10x + 25 = x^2 + 18x + 17 & We werken beide leden uit. De kwadraten vallen weg. \\
x = 1 & $8x = 8$, dus $x = 1$. \\
2,\ 6,\ 18 & We controleren door in te vullen: het quotiënt is telkens 3.
\end{stappen}}

\begin{oefeningen}
\oef{h2-k9}
\end{oefeningen}

\subsection{Grafiek}

Het voorschrift $u_n = u_1 \cdot q^{\,n-1}$ lijkt op dat van een exponentiële functie, met $n$ in plaats van $x$. Bij $q > 0$ liggen de punten dan ook op een exponentiële kromme. Het quotiënt bepaalt de vorm: bij $q > 1$ groeien de termen steeds sneller, bij $0 < q < 1$ krimpen ze naar 0 toe. Bij $q < 0$ wisselen de termen telkens van teken, en liggen ze om en om boven en onder de as. Zo'n rij hoort bij geen enkele exponentiële functie.

\figuren{\minigrafiek{$q = 2$}{-2}{34}{10,20,30}{1/1,2/2,3/4,4/8,5/16,6/32}\qquad
\minigrafiek{$q = \tfrac12$}{-2}{34}{10,20,30}{1/32,2/16,3/8,4/4,5/2,6/1}\qquad
\minigrafiek{$q = -\tfrac12$}{-10}{18}{-8,8,16}{1/16,2/-8,3/4,4/-2,5/1,6/-.5}}

\subsection{Exponentiële groei}

Meetkundige rijen beschrijven alles wat per tijdseenheid met hetzelfde \textit{percentage} groeit of krimpt: een bevolking, een spaarbedrag, een radioactieve stof. Waarom is dat een meetkundige rij? Kijk naar een eenvoudig geval: een bedrag van \eur 100 dat elk jaar met 10\,\% groeit.

\naast{Na een jaar komt er 10\,\% van 100 bij:
\[ 100 + 0{,}10 \cdot 100 = 100 \cdot (1 + 0{,}10) = 100 \cdot 1{,}1 = 110 . \]
Het jaar daarna komt er 10\,\% van 110 bij, en niet opnieuw 10 euro:
\[ 110 + 0{,}10 \cdot 110 = 110 \cdot 1{,}1 = 121 . \]
Telkens 10\,\% bijtellen is dus hetzelfde als telkens vermenigvuldigen met $1{,}1$. Na $n$ jaar heb je $n$ keer met $1{,}1$ vermenigvuldigd:
\[ u_n = 100 \cdot 1{,}1^{\,n} . \]}{\begin{tikzpicture}[x=.58cm,y=.022cm,font=\scriptsize]
  \draw[black!14,very thin,xstep=1,ystep=20] (0,0) grid (7.4,200);
  \draw[->,grijs] (0,0) -- (7.8,0) node[right,text=black] {$n$};
  \draw[->,grijs] (0,0) -- (0,212) node[above,text=black] {$u_n$};
  \foreach \t in {1,...,7}{\node[below,inner sep=1.5pt] at (\t,0) {$\t$};}
  \foreach \t in {40,80,...,200}{\node[left,inner sep=1.5pt] at (0,\t) {$\t$};}
  \foreach \x/\y in {0/100,1/110,2/120,3/130,4/140,5/150,6/160,7/170} \draw[grijs,thick] (\x,\y) circle (1.9pt);
  \foreach \x/\y in {0/100,1/110,2/121,3/133.1,4/146.41,5/161.05,6/177.16,7/194.87} \fill[accent] (\x,\y) circle (1.9pt);
  \node[right,text=accent] at (7.1,194.87) {$\cdot\, 1{,}1$};
  \node[right,text=grijs] at (7.1,165) {$+10$};
\end{tikzpicture}}

De figuur vergelijkt met lineaire groei, waarbij er elk jaar 10 euro bijkomt (de open bolletjes). Eerst is het verschil klein, maar het groeit snel: bij exponentiële groei levert ook de groei van vorig jaar zelf weer groei op.

Dezelfde redenering geldt voor elk percentage. Groeit een grootheid met $p\,\%$, dan komt er $\tfrac{p}{100}$ van het bedrag bij, en wordt het bedrag maal $1 + \tfrac{p}{100}$. Krimpt ze met $p\,\%$, dan gaat er $\tfrac{p}{100}$ af, en blijft er $1 - \tfrac{p}{100}$ van over. Die factor heet de \textit{groeifactor}. Omdat we starten bij de beginwaarde $u_0$, zijn er na $n$ tijdseenheden $n$ sprongen gezet.

\begin{theorie}
\emph{Exponentiële groei}: groeit een grootheid per tijdseenheid met $p\,\%$, dan is de \emph{groeifactor} $\formule{q = 1 + \dfrac{p}{100}}$. Bij een daling met $p\,\%$ is hij $\formule{q = 1 - \dfrac{p}{100}}$. Met beginwaarde $u_0$ is de grootheid na $n$ tijdseenheden

\vspace{1mm}
\qquad $\formule{u_n = u_0 \cdot q^{\,n}}$
\end{theorie}

Percentages tel je dus niet zomaar op. Tien jaar lang 2\,\% groei is geen 20\,\%, want $1{,}02^{10} \approx 1{,}219$: het is bijna 22\,\%. En wie 20\,\% verliest en daarna 20\,\% wint, staat niet terug op het begin: $0{,}8 \cdot 1{,}2 = 0{,}96$.

\voorbeeld{}{Een stad telt 8000 inwoners en groeit met 2\,\% per jaar. Hoeveel inwoners zijn er na 10 jaar? Wanneer zijn er 12\,000 inwoners?\label{vb:h2-mgroei}}{%
\uitleg{Een vast percentage per jaar: dat is exponentiële groei. De 8000 inwoners zijn er bij de start, dus $u_0 = 8000$. Voor de tweede vraag is de exponent onbekend, en hebben we een logaritme nodig.}
\begin{stappen}
q = 1 + \dfrac{2}{100} = 1{,}02 & Een groei van 2\,\% geeft groeifactor $1{,}02$. \\[1.5ex]
u_n = 8000 \cdot 1{,}02^{\,n} & Dit is het aantal inwoners na $n$ jaar. \\
u_{10} = 8000 \cdot 1{,}02^{10} \approx 9752 & Na 10 jaar zijn er ongeveer 9752 inwoners. \\
8000 \cdot 1{,}02^{\,n} = 12\,000 & We zoeken wanneer het aantal 12\,000 wordt. \\
1{,}02^{\,n} = 1{,}5 \ \Rightarrow \ n = \log_{1{,}02} 1{,}5 \approx 20{,}5 & We delen door 8000 en nemen de logaritme. Na 20 jaar zijn er nog geen 12\,000, na 21 jaar wel: in de loop van het 21ste jaar.
\end{stappen}}

\voorbeeld{}{Een nieuwe smartphone kost \eur 800 en verliest elk jaar 30\,\% van zijn waarde. Wat is hij na 4 jaar nog waard?\label{vb:h2-daling}}{%
\uitleg{Nu daalt de waarde met een vast percentage. Na elk jaar blijft er $70\,\%$ over, dus de groeifactor is kleiner dan 1.}
\begin{stappen}
q = 1 - \dfrac{30}{100} = 0{,}7 & Een daling van 30\,\% geeft groeifactor $0{,}7$. \\[1.5ex]
u_4 = 800 \cdot 0{,}7^{4} \approx 192{,}08 & Na 4 jaar is de smartphone nog ongeveer \eur 192 waard. Wie elk jaar \eur 240 aftrekt, zou op $800 - 4 \cdot 240 = -160$ uitkomen: dat kan niet, want er gaat telkens 30\,\% af van een kleiner bedrag.
\end{stappen}}

\nujij{Een stad telt 50\,000 inwoners en krimpt met 1,5\,\% per jaar. Hoeveel inwoners zijn er na 10 jaar?}{$u_{10} = 50\,000 \cdot 0{,}985^{10} \approx 42\,987$.}

\begin{oefeningen}
\oef{h2-k10}
\end{oefeningen}

%========================================================================================
\section{Enkelvoudige en samengestelde intrest}

Op een spaarrekening krijg je intrest: een percentage van je kapitaal. Er zijn twee manieren om die te berekenen. Bij \textit{enkelvoudige intrest} krijg je elk jaar intrest op het beginkapitaal alleen, dus elk jaar hetzelfde bedrag: dat is lineaire groei. Bij \textit{samengestelde intrest} komt de intrest bij het kapitaal, en krijg je het jaar daarna ook intrest op die intrest: dat is exponentiële groei.

\begin{theorie}
Je belegt een kapitaal $K_0$ aan $p\,\%$ per jaar. $K_n$ is het kapitaal na $n$ jaar.

\vspace{1mm}
{\renewcommand{\arraystretch}{1.3}%
\begin{tabular}{@{}>{\RaggedRight\arraybackslash}p{6.9cm}>{\RaggedRight\arraybackslash}p{6.9cm}@{}}
\textbf{Enkelvoudige intrest} & \textbf{Samengestelde intrest} \\
$\formule{K_n = K_0 \left(1 + n \cdot \dfrac{p}{100}\right)}$ & $\formule{K_n = K_0 \left(1 + \dfrac{p}{100}\right)^{n}}$
\end{tabular}}
\end{theorie}

De eerste formule is lineaire groei met $v = K_0 \cdot \tfrac{p}{100}$, de tweede exponentiële groei met $q = 1 + \tfrac{p}{100}$.

\voorbeeld{}{Je belegt \eur 2000 aan 3\,\% per jaar. Hoeveel heb je na 10 jaar, met enkelvoudige en met samengestelde intrest?\label{vb:h2-intrest}}{%
\uitleg{We vullen beide formules in met $K_0 = 2000$, $p = 3$ en $n = 10$, en vergelijken.}
\begin{stappen}
K_{10} = 2000 \,(1 + 10 \cdot 0{,}03) = 2600 & Enkelvoudig: elk jaar komt er \eur 60 bij, dus 10 keer \eur 60. \\
K_{10} = 2000 \cdot 1{,}03^{10} \approx 2687{,}83 & Samengesteld: \eur 87,83 meer, door de intrest op de intrest.
\end{stappen}
\naast{De figuur toont het kapitaal elk jaar: de open bolletjes bij enkelvoudige intrest, de volle bij samengestelde. De verticale as begint bij \eur 1800. In het begin is het verschil klein, maar het wordt elk jaar groter.}{\begin{tikzpicture}[x=.5cm,y=.0032cm,font=\scriptsize]
  \draw[black!14,very thin,xstep=1,ystep=200] (0,1800) grid (10.4,2800);
  \draw[->,grijs] (0,1800) -- (11,1800) node[right,text=black] {$n$};
  \draw[->,grijs] (0,1800) -- (0,2880) node[above,text=black] {$K_n$};
  \foreach \t in {2,4,...,10}{\node[below,inner sep=1.5pt] at (\t,1800) {$\t$};}
  \foreach \t in {2000,2400,2800}{\node[left,inner sep=1.5pt] at (0,\t) {$\t$};}
  \foreach \x in {0,...,10}{\draw[grijs,thick] (\x,{2000+60*\x}) circle (1.6pt);
    \fill[accent] (\x,{2000*1.03^\x}) circle (1.6pt);}
\end{tikzpicture}}}

\nujij{Je belegt \eur 1500 aan 6\,\% per jaar. Hoeveel heb je na 4 jaar, met enkelvoudige en met samengestelde intrest?}{Enkelvoudig \eur 1860; samengesteld ongeveer \eur 1893,72.}

\begin{oefeningen}
\oef{h2-k11}
\end{oefeningen}

%----------------------------------------------------------------------------------------
% OPDRACHTEN: het grootste deel van de oefeningen, op een nieuwe pagina
%----------------------------------------------------------------------------------------
\opdrachtenkop

\lkr{De opdrachten staan bewust door elkaar, met stijgende moeilijkheid. In deze versie staat bij elke opdracht het onderwerp, zodat je kunt kiezen. Geef ze pas volledig op als de hele leerstof gezien is, of selecteer per onderwerp.}

\begin{oefeningen}
\oef{h2-soort}

\oef{h2-o1}

\oef{h2-isterm}

\oef{h2-mvoorschrift}

\oef{h2-o2}

\oef{h2-tweetermen}

\oef{h2-o3}

\oef{h2-o4}

\oef{h2-misterm}

\oef{h2-voorschrift}

\oef{h2-mtwee}

\oef{h2-o5}

\oef{h2-o6}

\oef{h2-intrest}

\oef{h2-o7}

\oef{h2-o8}

\oef{h2-o9}

\oef{h2-groei}

\oef{h2-herken}

\oef{h2-o10}

\oef{h2-o11}

\oef{h2-o12}

\oef{h2-o13}

\oef{h2-o14}

\oef{h2-o15}

\oef{h2-o16}

\oef{h2-o17}

\oef{h2-o18}

\oef{h2-o19}

\oef{h2-o20}

\oef{h2-o21}

\oef{h2-o22}

\oef{h2-o23}

\oef{h2-o24}

\oef{h2-o25}

\oef{h2-o26}

\oef{h2-o27}

\oef{h2-o28}

\oef{h2-o29}

\oef{h2-o30}

\oef{h2-o31}

\oef{h2-o32}

\oef{h2-o33}

\oef{h2-o34}

\oef{h2-o35}

\oef{h2-o36}

\oef{h2-o37}

\oef{h2-o38}

\oef{h2-o39}

\oef{h2-o40}

\oef{h2-o41}

\oef{h2-o42}

\oef{h2-o43}

\oef{h2-o44}

\oef{h2-o45}

\oef{h2-o46}

\oef{h2-o47}

\oef{h2-o48}

\oef{h2-o49}

\oef{h2-o50}

\oef{h2-o51}

\oef{h2-o52}

\oef{h2-o53}

\oef{h2-o54}

\end{oefeningen}

%----------------------------------------------------------------------------------------
% KOPPEL: verbanden met andere leerstof. Enkel waar het zinvol is; geen koppeling tussen
% leerstof die vanzelf uit elkaar volgt (zoals limieten, afgeleiden en integralen).
%----------------------------------------------------------------------------------------
\opdrachtsectie{Koppel}

\lkr{Kijk na of deze leerstof al gezien is in jouw jaarplanning. Kansrekenen staat bv. ook in het 6de jaar.}

\begin{oefeningen}
\oef{h2-koppel-1}

\oef{h2-koppel-2}

\oef{h2-koppel-3}

\oef{h2-koppel-4}

\needspace{14\baselineskip}% figuur en deelvragen bij de opgave houden
\oef{h2-koppel-5}

\oef{h2-koppel-6}
\end{oefeningen}

%----------------------------------------------------------------------------------------
\opdrachtsectie{Test jezelf}

\begin{oefeningen}
\oef{h2-t1}

\oef{h2-t2}

\oef{h2-t3}

\oef{h2-t4}

\oef{h2-t5}

\oef{h2-t6}

\oef{h2-t7}

\oef{h2-t8}

\oef{h2-t9}
\end{oefeningen}

\herhaalschema{%
  \fzv{oef:h2-t1}{oef:h2-soort}%
  \fv{oef:h2-t2}{vb:h2-term}{oef:h2-voorschrift}%
  \fv{oef:h2-t3}{vb:h2-mterm}{oef:h2-mvoorschrift}%
  \fv{oef:h2-t4}{vb:h2-tweetermen}{oef:h2-tweetermen}%
  \fv{oef:h2-t5}{vb:h2-mtwee}{oef:h2-mtwee}%
  \fv{oef:h2-t6}{vb:h2-isterm}{oef:h2-isterm}%
  \fv{oef:h2-t7}{vb:h2-misterm}{oef:h2-misterm}%
  \fv{oef:h2-t8}{vb:h2-groei}{oef:h2-groei}%
  \fv{oef:h2-t9}{vb:h2-intrest}{oef:h2-intrest}}

\ifzelfstudie
\section*{Maak je eigen samenvatting}
Sluit dit boek en schrijf in je schrift: wat is een rekenkundige en wat is een meetkundige rij, welke voorschriften horen erbij, hoe vind je de rij uit twee termen, hoe test je of een getal een term is, wat is het verschil tussen enkelvoudige en samengestelde intrest, en welke fouten wil je vermijden? Vergelijk daarna met de kaders.
\fi
